\documentclass[12pt,a4paper]{article}
\usepackage[utf8]{inputenc}
\usepackage{geometry}
\geometry{left=3cm,right=3cm,top=3cm,bottom=3cm}
\usepackage{hyperref}
\begin{document}

\begin{flushright}
Yiwen Sun\\
School of Mathematics and Statistics\\
Beijing Technology and Business University\\
Beijing 100048, China\\
yver\_sun@163.com\\[1em]
Yanlin Shi (Corresponding Author)\\
Associate Professor\\
School of Languages and Communication\\
Beijing Technology and Business University\\
Beijing 100048, China\\
shiyanlin@th.btbu.edu.cn\\[1em]
Jiarui Liang\\
School of Mathematics and Statistics\\
Beijing Technology and Business University\\
Beijing 100048, China\\
2151460699@qq.com\\[1em]
\today
\end{flushright}

\vspace{1cm}

\noindent Dear Editor of \textit{Regional Studies},

\vspace{0.5cm}

I am writing to submit our manuscript entitled, ``\textbf{Policy Attention versus Institutional Supply: How Local Governments Foster Cluster Resilience---A Mixed-Methods Case Study of China's Gaoyang Textile Industry},'' for consideration as an Original Article in \textit{Regional Studies}.

The manuscript engages directly with the journal's core interest in the conceptual and empirical advancement of regional economic resilience. While the regional resilience literature has grown substantially since Martin (2012) and Boschma (2015), two gaps persist: institutional mechanisms are often treated as a black box in macro-level analyses, and empirical work is concentrated on advanced-economy clusters, leaving traditional clusters in developing economies understudied.

Our study addresses both gaps through three contributions:

\begin{enumerate}
    \item \textbf{Conceptual:} We introduce a distinction between \textit{policy attention} (discursive emphasis in government documents, measured via BERTopic dynamic topic modeling) and \textit{institutional supply} (actual resource allocation). This distinction travels beyond our case and provides a portable framework for any setting where textual policy signals proxy for government action.

    \item \textbf{Theoretical:} We articulate two sequential micro-mechanisms through which local government intervention can support cluster resilience---\textit{compliance cost socialization} via collective environmental infrastructure, and \textit{transaction cost reduction} via e-commerce and regional branding---grounded in a microeconomic formalization of firm cost structures. This extends evolutionary economic geography frameworks by specifying micro-level transmission channels.

    \item \textbf{Empirical:} Through a 25-year (2000--2024) mixed-methods case study of China's largest towel production cluster (Gaoyang County, Hebei Province), we trace a three-phase policy evolution and two-stage resilience trajectory using BERTopic topic modeling, Synthetic Control Method, and Interrupted Time Series analysis. We transparently disclose archival data constraints---including systematic heterogeneity in government report length between early (summary) and recent (full-text) periods---and explicitly characterize quantitative evidence as mechanism-illustrative rather than causal-confirmatory, given single-case limitations and limited statistical power.
\end{enumerate}

We believe this manuscript aligns with \textit{Regional Studies}' readership across economic geography, regional policy, and evolutionary economic geography. The conceptual framework, theoretical mechanisms, and transparent mixed-methods template are offered as contributions that can be tested, extended, and challenged in future comparative work. All data and code are openly available at \url{https://github.com/yver-sun/gaoyang-towel-resilience}.

This manuscript is original, has not been published previously, and is not under consideration for publication elsewhere. All authors have approved the manuscript and agree with its submission. The authors declare no conflict of interest.

\noindent\textbf{Suggested Reviewers:}
\begin{enumerate}
    \item Prof. Robert Hassink, Kiel University, Germany --- hassink@geographie.uni-kiel.de (regional resilience, evolutionary economic geography, East Asian old industrial areas)
    \item Prof. Xiaohui Hu, Nanjing Normal University, China --- huxh@njnu.edu.cn (economic resilience of old industrial cities in China, agency and institutions)
    \item Prof. Elvira Uyarra, University of Manchester, UK --- elvira.uyarra@manchester.ac.uk (regional innovation policy, policy mixes)
    \item Prof. Shengjun Zhu, Peking University, China --- shengjun.zhu@pku.edu.cn (industrial dynamics in China, environmental regulation, global production networks)
\end{enumerate}

Thank you for your time and consideration. We look forward to your response.

\vspace{1cm}

\noindent Sincerely,

\vspace{1cm}

\noindent Yiwen Sun\\
School of Mathematics and Statistics\\
Beijing Technology and Business University\\
Email: yver\_sun@163.com

\vspace{0.5cm}

\noindent Yanlin Shi (Corresponding Author)\\
Associate Professor\\
School of Languages and Communication\\
Beijing Technology and Business University\\
Email: shiyanlin@th.btbu.edu.cn

\vspace{0.5cm}

\noindent Jiarui Liang\\
School of Mathematics and Statistics\\
Beijing Technology and Business University\\
Email: 2151460699@qq.com

\end{document}
